\textbf{First Version}

Kepler third Law for satelite % sic: "satelite" in source
and Moon: \hfill(30 points)
\begin{align*}
  \frac{P_S^2}{a_S^3} &= \frac{P_S^2}{(15.3\,R_P)^3} = \frac{4\pi^2}{GM_P}; \\
  \frac{P_M^2}{a_M^3} &= \frac{P_M^2}{(60.3\,R_E)^3} = \frac{4\pi^2}{GM_E};
\end{align*}

Express mass in term of mass density: \hfill(30 points)
\begin{align*}
  \frac{P_S^2}{(15.3\,R_P)^3} &= \frac{4\pi^2}{GM_P}
    = \frac{4\pi^2}{G\rho_P \frac{4}{3}\pi R_P^3} \\
  \frac{P_M^2}{(60.3\,R_E)^3} &= \frac{4\pi^2}{GM_E}
    = \frac{4\pi^2}{G\rho_E \frac{4}{3}\pi R_E^3}
\end{align*}

Evaluate ratio of mass-density: \hfill(30 points)
\[ \frac{\rho_E}{\rho_P}
   = \frac{60.3^3 \times P_S^2}{15.3^3 \times P_M^2} \]
% sic: the source labels this ratio (and the next) rho_E/rho_P; the value
% 0.238 is actually rho_P/rho_E, as the Second Alternate Version confirms
\[ \frac{\rho_E}{\rho_P} = 0.238 \approx 0.24 \]
\hfill(10 points)

\textbf{Second Alternate Version}

\begin{itemize}
\item The period of the Satellite:
  $P_S = 7\ \text{days},\ 3\ \text{hours},\ 43\ \text{minutes}
   = 0.0196\ \text{years}$,
\item Mean orbital radius of the Satellite: $a_S = 15.3\,R_P$; $R_P$: radius
  of the Planet.
\item The period of the Moon:
  $P_M = 27\ \text{days},\ 7\ \text{hours},\ 43\ \text{minutes}
   = 0.0749\ \text{years}$,
\item Mean orbital radius of the Moon: $a_M = 60.3\,R_E$; $R_E$: radius of
  the Earth.
\item Mass of the Sun ($M$)
\item Kepler's Third Law: $G(M + m) = 4\pi^2 \dfrac{a^3}{T^2}$
  $\rightarrow$
  \begin{itemize}
  \item Mass of the Planet: $m_P$
  \item Radius of the Planet: $R_P$
  \item Semi major axis of the Planet $=$ Mean orbital radius of the
    Plant: % sic: "Plant" in source
    $a_P$
  \item Semi major axis of the Satellite $=$ Mean orbital radius of the
    Satellite: $a_S$
  \item Semi major axis of the Moon $=$ Mean orbital radius of the Moon:
    $a_M$
  \item Mass of the Earth: $m_E$
  \item Radius of the Earth: $R_E$
  \item Semi major axis of the Earth $=$ Mean orbital radius of the Earth:
    $a_E$
  \end{itemize}
\end{itemize}

\begin{description}
\item[Planet--Satellite:]
  $G(m_P + m_S) = 4\pi^2 \dfrac{a_S^3}{P_S^2}$, $m_P \gg m_S$
  $\rightarrow$ $G(m_P) = 4\pi^2 \dfrac{a_S^3}{P_S^2}$ \hfill(1)\quad(5 points)
\item[The Earth--the Moon:]
  $G(m_E + m_M) = 4\pi^2 \dfrac{a_M^3}{P_M^2}$, $m_E \gg m_M$
  $\rightarrow$ $G(m_E) = 4\pi^2 \dfrac{a_M^3}{P_M^2}$ \hfill(2)\quad(5 points)
\item[Planet--Sun:]
  $G(M + m_P) = 4\pi^2 \dfrac{a_P^3}{P_P^2}$, $M \gg m_P$
  $\rightarrow$ $G(M) = 4\pi^2 \dfrac{a_P^3}{P_P^2}$ \hfill(3)\quad(5 points)
\item[Earth--Sun:]
  $G(M + m_E) = 4\pi^2 \dfrac{a_E^3}{P_E^2}$, $M \gg m_E$
  $\rightarrow$ $G(M) = 4\pi^2 \dfrac{a_E^3}{P_E^2}$ \hfill(4)\quad(5 points)
\end{description}

From equation (1) and (3):
\[ \frac{m_P}{M}
   = \frac{4\pi^2 \frac{a_S^3}{P_S^2}}{4\pi^2 \frac{a_P^3}{P_P^2}}
   = \frac{(a_S)^3}{(a_P)^3}\,\frac{(P_P)^2}{(P_S)^2}
   = \frac{(15.3R_P)^3}{(a_P)^3}\,\frac{(P_P)^2}{(P_S)^2}
   = \frac{(15.3R_P)^3}{(a_P)^3}\,\frac{(P_P)^2}{(0.0196)^2} \]
\hfill(20 points)

From equation (2) and (4):
\[ \frac{m_E}{M}
   = \frac{4\pi^2 \frac{a_M^3}{P_M^2}}{4\pi^2 \frac{a_E^3}{P_E^2}}
   = \frac{(a_M)^3}{(a_E)^3}\,\frac{(P_E)^2}{(P_M)^2}
   = \frac{(60.3R_E)^3}{(a_E)^3}\,\frac{(P_E)^2}{(P_E)^2}
   = \frac{(60.3R_E)^3}{(a_E)^3}\,\frac{(P_E)^2}{(0.0749)^2}
   = \frac{(60.3R_E)^3}{(1)^3}\,\frac{(1)^2}{(0.0749)^2} \]
\hfill(20 points)

$\rightarrow$
\[ \frac{m_E}{m_P}
   = \frac{\dfrac{(15.3R_P)^3}{(a_P)^3}\,\dfrac{(P_P)^2}{(0.0196)^2}}
          {\dfrac{(60.3R_E)^3}{(1)^3}\,\dfrac{(1)^2}{(0.0749)^2}}
   = \frac{(15.3R_P)^3/(0.0196)^2}{(60.3R_E)^3/(0.0749)^2} \]
% sic: the source writes m_E/m_P for what is actually m_P/m_E
\hfill(20 points)

$\rightarrow$
\[ \frac{\rho_P}{\rho_E}
   = \frac{m_P/\mathrm{Vol}_P}{m_E/\mathrm{Vol}_E}
   = \frac{m_P/\frac{4}{3}\pi(R_P)^3}{m_E/\frac{4}{3}\pi(R_E)^3}
   = \frac{(15.3R_P)^3/(0.0196)^2}{(60.3R_E)^3/(0.0749)^2}
     \cdot \frac{\frac{4}{3}\pi(R_E)^3}{\frac{4}{3}\pi(R_P)^3}
   = \frac{\dfrac{(15.3)^3}{(0.0196)^2}}{\dfrac{(60.3)^3}{(0.0749)^2}}
   = 0.24 \]
\hfill(20 points)
